% Sections 10-17: taxes, earned income/GNI, transfers & disposable income, assets & capital
% formation, identities, sequence of accounts, MTA, past-paper answers.

%======================================================================
\section{Taxes and subsidies}\label{sec:tax}
\covers{--- (not on the list, but asked)}{Lecture 4}{2025--26 Q1(vii)}

Government has three economic roles: redistributing income through taxes and subsidies,
non-market production (``general government''), and market production through
quasi-corporations, which sit in the corporate sector.

\begin{itemize}
\item \df{Taxes} are compulsory, unrequited payments, in cash or kind, made to government units
  by households, enterprises and other institutional units.
\item \df{Subsidies} are current, unrequited payments by government \emph{to enterprises only},
  for production activity. They are based on the level of production: quantities or values of
  goods and services produced, sold or imported.
\item \df{Effects.} Taxes raise market prices and leave households and enterprises with fewer
  resources, which lowers private consumption and capital formation. Subsidies do the opposite.
\end{itemize}

\begin{keybox}{Where each kind of tax enters the accounts}
\small
\begin{tabularx}{\linewidth}{@{}l X@{}}
Taxes on products (GST, excise, import duties) & The wedge between basic and producers' prices.
Added to $\T{GVA}_{bp}$ to get GDP. Subsidies on products are the negative wedge. \\
Other taxes on production & Inside $\T{GVA}_{bp}$. Deducted before operating surplus. \\
Current taxes on income and wealth & \emph{Current transfers}. They reduce disposable income, not
GDP or GNI. \\
Capital taxes & \emph{Capital transfers}, in the capital account. \\
\end{tabularx}
\end{keybox}

%======================================================================
\section{Earned (primary) income and GNI}\label{sec:gni}
\covers{T1, T16, T17}{Lecture 8}{2025--26 Q1(i), Q2(ii)}

Household earnings are of two kinds. The first comes from taking part in production. The second
covers everything else: transfers, sale of property and holding gains. Only the first is
\df{earned (primary) income}.

\begin{itemize}
\item \df{Definition.} The income generated in production is the unit's GVA (at producers'
  prices). It is distributed to those who supplied factor services, as factor compensation, and
  to government, as production taxes less subsidies. The recipients' shares are their earned
  (primary) income.
\item \df{Who receives what.} Households and financial institutions receive CE, property income
  (interest and dividends on lent funds, rent on natural resources) and mixed income. Government
  receives taxes less subsidies on production and imports.
\item \df{Property income} is (i) investment income, meaning interest (adjusted for FISIM/IFSLD),
  distributed income of corporations such as dividends, and income on insurance and pension
  entitlements; plus (ii) \df{rent}, the payment for using \emph{non-produced} assets such as land
  or spectrum. A \df{rental} for \emph{produced} assets such as buildings and machinery is the
  purchase of a service, and so counts as IC.
\item Property income paid and received cancels within the domestic economy. That gives the
  \textbf{income approach}: $\T{GDP}=\T{RE}+\T{OS}+\T{MI}+(t-s)_{\text{products}}+(t-s)_{\text{imports}}$,
  with RE counting only what resident enterprises pay.
\item \df{Balance of primary income} $=$ primary income receivable $-$ primary income payable,
  worked out for a unit or a sector.
\item \df{GNI} is the sum of balances of primary income over all residents:
  GNI $=$ GDP $+$ net primary income (CE and property income) from the RoW. \textbf{GDP} measures
  production within the territory. \textbf{GNI} measures the income of residents.
\end{itemize}

\begin{tmpl}{Lecture 8 household $\to$ asked as 2025--26 Q2(ii)}
\small One member works in a factory and another runs a grocery. (a) Salary 5,000; (b) processed
food bought for resale 10,000; (c) remittance from abroad 5,000; (d) contribution to traders'
association 200; (e) income tax 500; (f) sales of processed food 15,000; (g) household provisions
4,350; (h) electricity for the grocery 150; (i) hired worker 1,000; (j) sales tax 150.\\
\textbf{In-flow} $=$ salary $+$ GVA of the grocery $=$ a $+$ (f $-$ b $-$ h) $=5{,}000+4{,}850=9{,}850$.
\textbf{Out-flow} $=$ i $+$ j $=1{,}000+150=1{,}150$. Balance of primary income $=8{,}700$.\\
c, d and e are transfers; g is final consumption. None of them is earned income.
\end{tmpl}

\begin{quick}{earned income? true of GNI?}
\small \textbf{Earned?} Beggars' income \no\ (transfer); smugglers' income \yes\ (illegal trade);
gamblers' income \no\ (transfer); gains from buying and selling shares \no\ (holding gain);
landlord's rent on land \yes; product taxes \yes; sale of land or an existing building \no; income
tax revenue \no\ (transfer).\\
\textbf{GNI?} Includes income tax received from the RoW \no\ (that is GNDI); measures production
\no\ (that is GDP); $=$ GDP $+$ net primary income from the RoW \yes; includes money from a
relative abroad \no; measures the earned income of residents \yes; includes receipts for
consultancy to the RoW \yes; unaffected by domestic production \no.
\end{quick}

%======================================================================
\section{Transfers, disposable income, consumption and saving}\label{sec:di}
\covers{T1, T10, T18, T19, T21}{Lecture 8}{2025--26 Q1(iv)}

\df{Transfers} are one-way: a unit provides goods, services, an asset or money and receives
nothing in return. Capital transfers are covered in \S\ref{sec:assets}.
\df{Current transfers}, other than social transfers in kind, are of three kinds:
\begin{enumerate}
\item current taxes on income, wealth, etc.;
\item \emph{net social contributions and social benefits}. Contributions are paid only by
  households, actual or imputed, and employers' contributions pass through CE first. They are
  received by government and by insurance and pension funds. Benefits are received by households
  when an event occurs: sickness, unemployment, retirement, housing, education, family;
\item \emph{other current transfers}: net non-life premiums and claims, transfers between
  government units (including foreign governments), transfers to and from NPISHs, and transfers
  between households, including remittances.
\end{enumerate}
Insurance premiums are \emph{partly} a payment for the insurer's service; claims are pure transfers.

\begin{quick}{current transfers?}
\small Entirely: corporate tax; election bonds; money sent home by a non-resident; insurance claims
received by a company; a worker's union contribution. Partly: life and non-life premiums. Not
current: a donation to build a school (capital transfer); royalty for a patented design (payment
for a service).
\end{quick}

\begin{keybox}{Disposable income}
\small \textbf{Disposable income} $=$ balance of primary income $+$ current transfers receivable
$-$ current transfers payable, leaving out social transfers in kind. It measures purchasing power.\\
\textbf{GNDI} $=$ GNI $+$ net current transfers from the RoW, including net taxes on income and
wealth $=$ GDP $+$ net primary income from the RoW $+$ net current transfers from the RoW.
\end{keybox}

\begin{tmpl}{Lecture 8 household, continued: the likely next variant}
\small Transfers in $=$ c $=5{,}000$; transfers out $=$ d $+$ e $=700$; net $=4{,}300$.
Disposable income $=8{,}700+4{,}300=13{,}000$. Extension: saving $=$ DI $-$ FCE (g)
$=13{,}000-4{,}350=8{,}650$.\\
\textbf{Symbolic} (2025--26 letters, with $S'$ for the grocery's sales, since the paper uses $S$
twice): $\text{DI}=S+S'-P-\sym{PE}-\sym{PHW}-\sym{PS}+R-\sym{PT}-\sym{PI}$, and saving
$=\text{DI}-\sym{PH}$.
\end{tmpl}

\df{Consumption} is using goods and services up. It takes three forms: \emph{completely} in
production (IC), \emph{partially} in production (CFC, i.e.\ depreciation), or for the direct
satisfaction of needs (\df{final consumption}). Final consumption meets \emph{individual} needs
(a household's own wants) or \emph{collective} needs (services to the whole community). Only
general government and the central bank provide collective services.

\begin{keybox}{HFCE and GFCE (asked: 2025--26 Q1(iv))}
\small
\textbf{Who has FCE:} households, general government, the central bank and NPISHs, but
\emph{never corporations}. $\T{FCE}=\T{PFCE}+\T{GFCE}+\T{CBCE}$, where $\T{PFCE}=\T{HFCE}+$ FCE of
NPISHs.\\[2pt]
\textbf{HFCE}: consumption spending by \emph{resident} households, both in the territory (on
domestic output and imports) and abroad (counted as imports). It includes own-account consumption
such as crops kept and owner-occupied housing. Spending by non-residents in the territory is
exports, not HFCE.\\[2pt]
\textbf{GFCE}: the whole non-market output of general government is treated as consumed by
government itself, even though households enjoy it. Add purchases of market goods passed free to
households, e.g.\ food for flood victims or blankets for the homeless:\\[1pt]
\centerline{$\T{GFCE}=\text{total output of GG}-\text{own-account GCF}-\text{sales of goods and services}
+\text{benefits in kind}.$}
\end{keybox}

\begin{tmpl}{Lecture 8 Malaysian farmer}
\small Mixed income 5,000; rent received 1,000. Grain kept: 100 for seed, 900 to eat. Consumption
goods bought 4,000; fertiliser 500. $\T{FCE}=900+4{,}000=4{,}900$. The seed is inventory, i.e.
next year's IC, and fertiliser is IC.\\
With current transfers of 2,000 and CFC 500: DI $=5{,}000+1{,}000+2{,}000=8{,}000$; gross saving
$=8{,}000-4{,}900=3{,}100$; net saving $=3{,}100-500=2{,}600$.\\
\textbf{Saving:} gross saving $=$ disposable income $-$ FCE (GNDI $-$ FCE for the economy);
net saving $=$ gross saving $-$ D\&D.
\end{tmpl}

%======================================================================
\section{Assets, capital formation and net lending}\label{sec:assets}
\covers{T8, T20, T21, T22}{Lecture 9}{2025--26 Q1(v)}

A \df{financial asset} comes from a contract between two units that gives one of them the right to
a payment, or a series of payments, from the other. Currency, deposits, loans, shares and bonds
are examples. A \df{liability} is the debtor's obligation to pay the creditor. All liabilities are
financial. The payment due is the financial claim, and it ends when the debtor pays.

\begin{keybox}{Non-financial assets (2025 SNA) --- produced vs non-produced (asked: 2025--26 Q1(v))}
\small
\textbf{Produced assets} come into existence as outputs of production:
\begin{itemize}
\item \emph{fixed assets}, used repeatedly in production for more than a year: dwellings, other
  buildings and structures, machinery and equipment, weapons systems, costs of ownership transfer
  on non-produced assets, and IPP;
\item \emph{inventories}: materials and supplies, work in progress, finished goods, military
  inventories, goods for resale;
\item \emph{valuables}, held as stores of value: precious metals and stones, antiques and art,
  jewellery. They are neither consumed nor used in production.
\end{itemize}
\textbf{Non-produced assets} come into existence other than by production: contracts, leases and
licences; purchased goodwill and marketing assets.\\
\textbf{Natural resources}: land, minerals, energy resources (renewable and not), water,
non-cultivated biological resources (non-produced), cultivated biological resources (produced),
and radio spectra (either).\\[2pt]
\textbf{Treatment differs.} Produced fixed assets \emph{depreciate}, and acquiring them is GCF.
Natural resources \emph{deplete}, and acquiring them is recorded separately, not in GCF. The
2008 SNA introduced land improvements, weapons systems and IPP as fixed assets.
\end{keybox}

\begin{keybox}{Capital formation (T21, T22)}
\small
$\T{GDCF}\ (=\T{GCF}) = \T{GFCF} + \T{CII} + \text{acquisitions less disposals of valuables}$.
\begin{itemize}
\item \textbf{GFCF}: acquisitions less disposals of produced fixed assets, whether new or
  existing, and whether bought, bartered, received as capital transfers or made on own account;
  \emph{plus} costs of ownership transfer on non-produced assets; \emph{plus} major improvements
  that extend asset lives. Only GFCF is ``investment'' in the narrow sense.
\item \textbf{CII} $=$ closing stock $-$ opening stock of inventories: unsold finished goods,
  traders' stocks, unused raw materials and work in progress.
\end{itemize}
\end{keybox}

\df{Capital transfers} transfer ownership of a fixed asset, are tied to acquiring or disposing
of one, or forgive a liability. Examples are grants to build a school or temple, debt waivers and
capital taxes. Gross saving plus net capital transfers is usually less than GCF, and the gap is
borrowed:
\begin{center}
Net lending ($+$) / borrowing ($-$) $=$ gross saving $+$ net capital transfers $-$ GCF\\
$-$ net acquisition of non-produced non-financial assets.
\end{center}

\begin{tmpl}{Lecture 9 NPI}
\small A free-education NPI receives donations of 10,000, pays salaries of 7,000, pays no tax, and
has CFC of 200. It gets a capital grant of 50,000 for a classroom costing 100,000.\\
Saving $=10{,}000-7{,}000-200=2{,}800$; funds available $=2{,}800+50{,}000=52{,}800$; it must
borrow $100{,}000-52{,}800=47{,}200$.\\
\textbf{Symbolic.} Donation $D$, salaries $W$, CFC $C$, grant $G$, cost $K$: saving $=D-W-C$;
borrowing $=K-(D-W-C+G)$.
\end{tmpl}
\begin{trap}{CFC is not cash}
\small The slides' 47,200 treats CFC as reducing the funds. Strictly, gross saving is 3,000, so
SNA net borrowing $=100{,}000-3{,}000-50{,}000=47{,}000$. In the exam, follow the slides' working;
if you have time, add a line noting that CFC is non-cash.
\end{trap}

%======================================================================
\section{The main identities}\label{sec:id}
\covers{T3, T23}{Lectures 2, 9}{---}

\begin{keybox}{Three approaches to GDP}
\vspace{-8pt}
\begin{align*}
\text{Production:}\quad \T{GDP} &= \T{GVO}_{bp}-\T{IC}+(t-s)_{\text{products}}+(t-s)_{\text{imports}}\\
\text{Income:}\quad \T{GDP} &= \T{RE}+\T{OS}+\T{MI}+(t-s)_{\text{products}}+(t-s)_{\text{imports}}\\
\text{Expenditure:}\quad \T{GDP} &= \T{PFCE}+\T{GFCE}+\T{CBCE}+\T{GFCF}+\T{CII}+\text{valuables}+X-M
\end{align*}
\end{keybox}

\textbf{The eight numbered identities (Lecture 9)}, which underlie the sequence of accounts:
\begin{enumerate}[label={[\arabic*]},leftmargin=2.2em]
\item \emph{Commodity balance}: $\T{GVO}+M \equiv \T{IC}+\T{PFCE}+\T{GFCE}+\T{CBCE}+\T{GFCF}+\T{CII}+\text{valuables}+X$.
  Everything produced or imported is used.
\item \emph{Production}: $\T{GDP}\equiv\T{GVO}_{bp}-\T{IC}+(t-s)_{\text{products}}+(t-s)_{\text{imports}}$.
\item \emph{Income}: $\T{GDP}\equiv(\T{RE}+\T{OS}+\T{MI})$ generated in domestic enterprises
  $+\,(t-s)_{\text{products}}+(t-s)_{\text{imports}}$.
\item \emph{National income}: GNI $\equiv$ [3] $+$ net RE from the RoW $+$ net property income
  from the RoW.
\item \emph{Disposable income}: GNDI $\equiv$ GNI $+$ net current transfers $+$ net taxes on
  income and wealth from the RoW.
\item \emph{Expenditure}: $\T{GDP}\equiv\T{PFCE}+\T{GFCE}+\T{CBCE}+\T{GFCF}+\T{CII}+\text{valuables}+X-M$.
\item \emph{Saving}: $\text{Gross saving}\equiv\T{GNDI}-\T{PFCE}-\T{GFCE}-\T{CBCE}$.
\item \emph{Net lending}: $\equiv$ gross saving $+$ net capital transfers $-(\T{GFCF}+\T{CII}+\text{valuables})$
  $-$ net acquisition of non-produced non-financial assets.
\end{enumerate}

\df{Commodity-flow identity} (Lecture 2), with both sides at purchasers' prices:
\[
\T{GVA}_{bp}+(t-s)_{\text{products}}+M+\text{import duties}=\T{HFCE}+\T{GCF}+X,
\qquad\text{i.e.}\qquad Y+M=C+I+X.
\]

\begin{keybox}{The ladder of aggregates}
\small
\begin{tabularx}{\linewidth}{@{}r X r@{}}
& GDP & \\
$+$ & net primary (earned) income from RoW & $=$ \textbf{GNI} \\
$+$ & net current transfers (incl.\ taxes on income and wealth) from RoW & $=$ \textbf{GNDI} \\
$-$ & final consumption expenditure & $=$ \textbf{Gross saving} \\
$+$ & net capital transfers from RoW $-$ GCF $-$ net acquisition of non-produced assets &
$=$ \textbf{Net lending/borrowing} \\
\end{tabularx}\\[2pt]
Subtract D\&D from any of GDP, GNI, GNDI or saving to get the net version.
\end{keybox}

\begin{trap}{D\&D in the last rung}
\small The Lecture 9 summary slide also subtracts D\&D in the final step. That is only right if
you start from \emph{net} saving and use \emph{net} capital formation (GCF $-$ D\&D). Starting
from gross saving, do \emph{not} subtract D\&D. The Lecture 12 capital account confirms this:
gross saving $33+10=43$, and $43+1-2-23=19$.
\end{trap}

%======================================================================
\section{The sequence of accounts and balance sheets}\label{sec:seq}
\covers{T24}{Lecture 10}{---}

The SNA records all flows and stocks. Accounts are compiled for each of the five sectors, the RoW
and the total economy. They are consistent because there is one production boundary and one asset
boundary. There are three classes:
\begin{enumerate}
\item \df{Current accounts} record production, income and its use. They do not change the balance
  sheet directly. They are the \emph{production account} and the \emph{income accounts}:
  generation of income, allocation of primary income, secondary distribution (income transfers
  other than social transfers in kind), and use of disposable income.
\item \df{Accumulation accounts} record changes in assets and liabilities: the \emph{capital
  account} (capital formation, capital transfers, non-produced assets) and the \emph{financial
  account} (all transactions in financial assets and liabilities). Other changes in volume and
  revaluation also belong here but are not in this course.
\item \df{Balance sheets} record stocks at the opening and closing of the period. Assets
  (produced, non-produced, financial) sit on the left; liabilities (all financial) and net worth
  on the right. \textbf{Net worth} $=$ total assets $-$ total liabilities, the measure of wealth.
  Closing balance sheet $=$ opening balance sheet $+$ transactions $+$ other flows.
\end{enumerate}

\begin{center}\small
\begin{tabular}{@{}lll@{}}
\toprule
\textbf{Account type} & \textbf{Left side} & \textbf{Right side} \\
\midrule
Current accounts & Uses (``expenditures'') & Resources (``revenues'') \\
Accumulation accounts & Changes in assets & Changes in liabilities and net worth \\
Balance sheet & Assets & Liabilities and net worth \\
\bottomrule
\end{tabular}
\end{center}

Every account is a T-account under double entry, so its two sides total the same. The
\df{balancing item} is resources minus uses, shown on the uses side. It is carried forward as the
first resource of the next account. The balancing items are the headline aggregates: GDP/NDP,
OS/MI, GNI/NNI, disposable income, saving, and net lending/borrowing. The full worked chain is in
\S\ref{sec:mta}.

%======================================================================
\section{Monetary transactions and the MTA}\label{sec:mta}
\covers{T25, T26}{Lectures 11--12}{---}

\textbf{The compiler's view.} Two-party transactions can usually be observed. Intra-unit
transactions, some non-produced items, and \emph{composite} financial transactions (one payment
that bundles several analytically different transactions) cannot, so they are imputed or derived.

\begin{multicols}{2}\small
\df{Monetary transactions}: money changes hands, always between two units, and they are recorded
as observed. The exception is composite ones such as nominal interest and total insurance
premiums, which are split into components (FISIM).
\begin{itemize}
\item Goods and services, produced or non-produced (g\&s $\leftrightarrow$ money);
\item distributive: earned income (factor services $\leftrightarrow$ money) and transfers
  (money $\to$ nothing);
\item financial (money $\leftrightarrow$ money), simple or composite.
\end{itemize}
\columnbreak
\df{Non-monetary transactions}: no money flows, so they cannot be observed. They are imputed at
market-equivalent prices, or at the sum of costs if there is no market.
\begin{itemize}
\item \emph{Two-party}: barter and payments in kind; transfers in kind;
\item \emph{intra-unit}: production for own final use (own-account capital formation, own
  consumption), and D\&D.
\end{itemize}
\end{multicols}

\begin{keybox}{Money balance identity}
\centering money in-flow $\equiv$ money out-flow $+$ change in cash in hand\\[1pt]
\raggedright\small Wallet: starts with 100; lunch 20, transport 30, client pays 150, loan from a
friend 350, stationery 500. In-flow 500, out-flow 550, change $-50$, leaving 50.
\end{keybox}

\begin{multicols}{2}\small
\textbf{In-flows} come from selling products, receiving factor and property income, receiving
transfers, \emph{incurring liabilities} (taking a loan, issuing shares or bonds), and
\emph{disposing of financial assets}.\\
\columnbreak
\textbf{Out-flows} go to buying products, paying income, paying transfers, \emph{extinguishing
liabilities} (repaying a loan, buying back shares), and \emph{acquiring financial assets}
(deposits, shares, bonds). Cash is a financial asset.
\end{multicols}

\begin{tmpl}{Lecture 11 household MTA}
\small Sold produce 240, bought goods and services 200, received wages 20, paid wages 80, received
remittances 10, paid transfers 2, borrowed 100. In-flow 370, out-flow 282, so financial assets
rose by 88.\\[2pt]
\centering
\begin{tabular}{@{}rl|lr@{}}
\toprule
\multicolumn{2}{c|}{\textbf{Expenditures}} & \multicolumn{2}{c}{\textbf{Revenues}} \\
\midrule
200 & Bought products & Sold products & 240 \\
80 & Paid income & Received income & 20 \\
2 & Paid transfers & Received transfers & 10 \\
88 & Change in financial assets (incl.\ cash) & Change in liabilities & 100 \\
\midrule
370 & Total & Total & 370 \\
\bottomrule
\end{tabular}
\end{tmpl}

\begin{itemize}
\item \df{Double entry.} Every monetary transaction affects two cells of a unit's MTA. A tractor
  (40) bought with a bank loan is ``bought products 40'' and ``change in liabilities 40''. Wages
  (10) paid into a bank account are ``received income 10'' and ``change in financial assets 10''.
  Each transaction has two parties, so the system as a whole records it four times
  (\df{quadruple entry}).
\item \df{Aggregation} adds corresponding cells. Among residents only, products sold $\equiv$
  bought, income receivable $\equiv$ payable, transfers receivable $\equiv$ payable, and change
  in liabilities $\equiv$ change in financial assets.
\item \df{With the RoW}, the two sides no longer match item by item:
  products sold $+$ imports $\equiv$ products bought $+$ exports, and so on. The one identity
  that holds for \emph{every} sector and for the RoW is
  \begin{center}
  sub-total non-financial (revenue) $+$ change in liabilities\\
  $\equiv$ sub-total non-financial (expenditure) $+$ change in financial assets,
  \end{center}
  so \textbf{non-financial balance $=$ net lending $=$ change in financial assets $-$ change in
  liabilities}.
\end{itemize}

\begin{center}\small
\begin{tabular}{@{}rl|lr|r@{}}
\toprule
\multicolumn{2}{c|}{\textbf{Expenditure}} & \multicolumn{2}{c|}{\textbf{Revenue}} & \textbf{Rev $-$ Exp} \\
\midrule
1235 & Bought products & Sold products & 1259 & 24 \ ($=X-M$) \\
404 & Paid income & Received income & 400 & $-4$ \\
5 & Paid transfers & Received transfers & 4 & $-1$ \\
1644 & \textit{Sub-total non-financial} & \textit{Sub-total non-financial} & 1663 & \textbf{19} \\
203 & Change in financial assets & Change in liabilities & 184 & $-19$ \\
\bottomrule
\end{tabular}
\end{center}
This is the aggregated MTA of all residents (Lecture 11), and it shows the economy is a
\textbf{net lender of 19}. The same identity holds sector by sector in the Lecture 12 table.
Households: $405+4=401+8$, so they lend 4. RoW: $323-342=-19$, so the RoW borrows the 19.

\df{From the MTA to the SNA accounts} (Lecture 12). The MTA records only monetary transactions and
does not say what purchases were for. Extending it takes three steps: split it by sector; add the
non-monetary transactions (own-use output 23, non-market output 43, so $\T{GVO}_{mp}=1259+23+43=1325$);
and break the broad items down. Use of products, 1301, splits into IC 870, FC 408 and CF 23.
Income splits into RE, property income and $(t-s)$. Transfers split into current and capital. With
D\&D $=10$:

\begin{center}\small
\begin{tabularx}{\textwidth}{@{}l X l@{}}
\toprule
\textbf{Account} & \textbf{Resources $-$ uses} & \textbf{Balancing item} \\
\midrule
\multicolumn{3}{@{}l}{\textit{Current accounts}} \\
Production & GVO$_{mp}$ 1325 $-$ IC 870 & GDP 455; NDP 445 \\
Generation of income & NDP 445 $-$ RE paid 153 $-$ production $(t-s)$ and import duties 130
  & OS $+$ MI (net) 162 \\
Allocation of primary income & 162 $+$ RE received 152 $+$ PI received 118 $+$ $(t-s)$ 130
  $-$ PI paid 121 & NNI 441 \\
Secondary distribution & 441 $+$ current transfers received 3 $-$ current transfers paid 3
  & NNDI 441 \\
Use of disposable income & 441 $-$ final consumption 408 & net saving 33 \\
\midrule
\multicolumn{3}{@{}l}{\textit{Accumulation accounts}} \\
Capital & 33 $+$ capital transfers received 1 $-$ paid 2 $-$ (GCF 23 $-$ D\&D 10)
  & net lending 19 \\
Financial & change in financial assets 203 $-$ change in liabilities 184 & net lending 19 \\
\bottomrule
\end{tabularx}
\end{center}
Check by expenditure: $408+23+24=455=$ GDP. Net lending is the \emph{same} 19 as the MTA balance,
because each non-monetary item enters both sides.

\begin{keybox}{Rules of accounting (Lecture 12)}
\small \textbf{Valuation.} Monetary exchanges are valued at the price the buyer actually paid.
Non-monetary ones use market-price equivalents, or the sum of costs where there is no market.
Exports and imports are f.o.b. Flows and stocks are both at current market value.\\
\textbf{Timing.} Products and financial assets are recorded when economic ownership changes.
Primary income and transfers are recorded when they fall due (accrual). Output and IC are
recorded when production takes place. Both parties must record the same value in the same period.
\end{keybox}

\begin{quick}{Lectures 11--12 true/false}
\small Cash in hand is not a financial asset \no. Every monetary transaction makes at least two
entries in a unit's MTA \yes. The MTA records stocks of assets \no. Unequal MTA totals mean a
missing entry \yes. Borrowing is a money in-flow \yes. In an open economy, products bought equal
products sold \no. The non-financial balance is net lending \yes. Money lent between residents
cancels out \yes. The MTA records detailed production, consumption and accumulation \no. All
transaction identities fail for sector accounts \no. The RoW account is presented from the
non-residents' point of view \yes. Net lending is affected by non-monetary transactions \no. Crop
output for own use is recorded in the use-of-income account \yes. Import duties are in the
production account \yes. All non-monetary transactions are valued at cost \no. Primary income and
transfers are on an accrual basis \yes.
\end{quick}

%======================================================================
\section{Past papers: model answers}\label{sec:pyq}

\subsection*{2025--26 Q1 --- define and relate, any 5 of 8 \hfill [(2.5 + 2.5) $\times$ 5 = 25]}
Give all eight a look, then pick the five you can write fastest. The shape of each answer:
define the first term, define the second, and give the relation, preferably as an equation.

\textbf{(i) GNI and GDP.} \emph{GDP} is the value of production by resident producers in the
economic territory during the period: $\sum\T{GVA}_{bp}+(t-s)_{\text{products}}+(t-s)_{\text{imports}}$.
It can also be computed as $\T{RE}+\T{OS}+\T{MI}+(t-s)$, or as $\T{PFCE}+\T{GFCE}+\T{CBCE}+\T{GCF}+X-M$.
\emph{GNI} is the sum of the balances of primary income of all residents, i.e.\ their earned
income. The relation is $\T{GNI}=\T{GDP}+$ (CE $+$ property income) received from the RoW $-$
(CE $+$ property income) paid to the RoW. GDP is a \emph{territorial production} concept and GNI
a \emph{residents' income} concept. Remittances and income taxes from abroad are transfers, so
they enter GNDI, not GNI. $\T{NNI}=\T{GNI}-\T{D\&D}$. [\S\ref{sec:prod}, \S\ref{sec:gni}]

\textbf{(ii) GVA and GVO.} \emph{GVO} is the value of the goods and services an establishment
produces in the period. Market output is valued at basic prices (quantity $\times$ price);
non-market and own-use output at the sum of costs. It includes changes in inventories of finished
goods and work in progress, and own-account capital formation. \emph{GVA} $=$ GVO $-$ IC, the
value the unit itself adds. It is shared out as CE $+$ OS/MI $+$ other net taxes on production.
Summing $\T{GVA}_{bp}$ over residents and adding net product taxes gives GDP; NVA $=$ GVA $-$
D\&D. [\S\ref{sec:prod}, \S\ref{sec:os}]

\textbf{(iii) Basic, producers' and purchasers' prices.} Define all three (\S\ref{sec:prices}) and
give the chain. Purchasers' price $-$ TTM $=$ producers' price. Producers' price $-$ (taxes $-$
subsidies) on products $=$ basic price. Output is valued at basic prices and uses at purchasers'
prices. Illustrate with the pen: 20, then 12, then 10.

\textbf{(iv) HFCE and GFCE.} Use the definitions in the box in \S\ref{sec:di}. The relation:
both are parts of FCE, with $\T{FCE}=\T{PFCE}+\T{GFCE}+\T{CBCE}$, and they enter
$\T{GDP}=\T{PFCE}+\T{GFCE}+\T{CBCE}+\T{GCF}+X-M$. Households \emph{consume} most government
output, but the \emph{expenditure} is recorded as government's. Corporations never have FCE.

\textbf{(v) Produced and non-produced assets.} Produced assets arise from production: fixed
assets, inventories, valuables. Non-produced assets arise otherwise: natural resources such as
land, minerals and water; contracts, leases and licences; goodwill. The treatment differs.
Produced fixed assets depreciate and their acquisition is GCF. Non-produced natural resources
deplete and their acquisition is a separate capital-account item. [\S\ref{sec:assets}]

\textbf{(vi) Depletion and depreciation.} Use the table in \S\ref{sec:dd}: which asset each
applies to, what causes it, current-price valuation, and that both are costs of production.
Then state $\text{Gross}-\text{Net}=\T{D\&D}$.

\textbf{(vii) Taxes and subsidies.} Give both definitions from \S\ref{sec:tax}: compulsory and
unrequited, versus current unrequited payments to enterprises only, based on the level of
production. Then the relation: product $(t-s)$ is the gap between producers' and basic prices, and
$\T{GDP}=\T{GVA}_{bp}+(t-s)_{\text{products}}$. Income taxes are current transfers, and taxes
raise prices while subsidies lower them.

\textbf{(viii) Transactions of current and capital nature.} Use the box in \S\ref{sec:flows}.
Current transactions involve production, primary income, final consumption and current transfers,
and go in the current accounts. Capital transactions involve capital formation, non-produced
assets, capital transfers and capital taxes, go in the accumulation accounts, and change the
balance sheet.

\subsection*{2025--26 Q2 --- answer one \hfill [5]}

\textbf{(i) Textile unit.} Garments $X_1$ (without taxes), fibres $X_2$, salaries $X_3$,
electricity and other incidentals $X_4$. [1 $+$ 2 $+$ 2]
\begin{align*}
\T{GVO}_{bp} &= X_1 \quad \text{(``without taxes'' means basic prices)}\\
\T{IC} &= X_2+X_4\\
\T{OS} &= \T{GVA}_{bp}-\T{CE}=(X_1-X_2-X_4)-X_3=X_1-X_2-X_3-X_4
\end{align*}
State the assumptions: no other taxes on production, and OS is gross, i.e.\ before depreciation.
Salaries are a factor payment, not IC.

\textbf{(ii) Household.} The paper uses $S$ both for salary (a) and for sales (f). Call the sales
$S'$ and \emph{say so in your answer}. [3 $+$ 2]
\begin{align*}
\text{In-flow} &= S+(S'-P-\sym{PE}) && (\text{salary} + \text{GVA of the grocery})\\
\text{Out-flow} &= \sym{PHW}+\sym{PS} && (\text{CE paid} + \text{product tax})
\end{align*}
Excluded: $R$, $\sym{PT}$ and $\sym{PI}$ are current transfers, and $\sym{PH}$ is final
consumption. The balance of primary income is $S+S'-P-\sym{PE}-\sym{PHW}-\sym{PS}$. The slide
numbers give 9,850, 1,150 and 8,700.
